\chapter{Literature Review}
\label{ch:literature_review}

The study of transverse momentum ($p_T$) spectra in high-energy proton-proton and heavy-ion collisions has a rich history, with phenomenological models continually evolving to capture both the soft (thermal) and hard (perturbative QCD) regimes of particle production. This chapter reviews the foundational models and recent advancements that form the baseline for our validation framework.

\section{Thermal and Hydrodynamical Foundations}
Early approaches to modeling particle spectra relied heavily on the Boltzmann-Gibbs statistics, assuming local thermal equilibrium of the produced fireball. However, the observation of power-law tails at high $p_T$ necessitated refinements. The Blast-Wave model \citep{Schnedermann1993} introduced collective transverse flow, successfully describing the soft bulk of the spectra ($p_T \lesssim 2.5$ GeV/c) by assuming a locally thermalized source expanding with a collective transverse velocity profile. 

The standard Boltzmann-Gibbs Blast-Wave (BGBW) model evaluates the invariant yield as an integral over the freeze-out surface. While highly successful at low $p_T$, the BGBW model systematically underpredicts the yield at higher transverse momenta, where hard parton scattering dominates.

\section{The Tsallis Distribution and Non-Extensive Thermodynamics}
To simultaneously describe both the low-$p_T$ thermal bulk and the high-$p_T$ power-law tail, the non-extensive Tsallis statistics were introduced to high-energy physics. A seminal advancement was the formulation of a thermodynamically consistent Tsallis distribution by \citet{Cleymans2012}, which strictly satisfies thermodynamic relations for particle number, energy density, and pressure. The Tsallis distribution introduces a non-extensivity parameter $q$; as $q \to 1$, the distribution recovers the standard Boltzmann-Gibbs form.

\section{Recent Advancements and Multiplicity Dependence}
Primary measurements of charged-particle $p_T$ spectra at the LHC have been extensively provided by the ALICE Collaboration \citep{ALICE2018}. These foundational datasets spanning multiple collision systems ($pp$, $p$-Pb, and Pb-Pb) and collision energies (up to $\sqrt{s} = 13$ TeV) serve as the empirical bedrock for phenomenological modeling.

Recent studies have focused heavily on the multiplicity dependence of freeze-out parameters in these proton-proton collisions. \citet{Khuntia2019} analyzed $p_T$ spectra at $\sqrt{s} = 7$ TeV using both the BGBW and the thermodynamically consistent Tsallis distribution. They observed a mass ordering in the kinetic freeze-out temperature $T_{\text{kin}}$ and the non-extensivity parameter, supporting a differential freeze-out scenario. Crucially, their BGBW analysis showed that the average radial flow velocity $\langle \beta \rangle$ remains largely independent of the event multiplicity, providing a key benchmark for multi-component models.

\citet{Rath2020} extended this work to 5.02 and 13 TeV collision energies. Their findings corroborated that $\langle \beta \rangle$ is almost independent of collision energy and multiplicity. However, they found that $T_{\text{kin}}$ significantly depends on the multiplicity classes. They concluded that while the Tsallis model better describes the full $p_T$ range, the BGBW model remains highly effective for the bulk part ($p_T \lesssim 2.5$ GeV/c).

More recently, \citet{Biro2025} proposed a two-component soft/hard baseline model for ALICE $pp$ data from 2.76 to 13 TeV. Their analysis revealed that a simple Boltzmann fit effectively describes the soft component across all energies, while the residual hard spectra exhibit no evolution in shape or peak position with event multiplicity. Furthermore, the mean $p_T$ for both the soft and hard components remains nearly constant across multiplicity classes, challenging more complex hydrodynamical interpretations that necessitate highly parameterized models.

These literature benchmarks underscore the necessity of rigorous statistical validation when proposing novel multi-component models, such as the 3-component formulation (incorporating a moving J\"uttner distribution) evaluated in this thesis.

To systematically contextualize the evolution of phenomenological modeling in this domain, Table \ref{tab:pt_models} provides a summary of seminal and contemporary works applying various statistical functions to $p_T$ spectra.

\begin{table}[h]
\centering
\caption{Summary of phenomenological models applied to transverse momentum spectra in high-energy collisions.}
\label{tab:pt_models}
\resizebox{\textwidth}{!}{
\begin{tabular}{|l|l|l|l|l|l|}
\hline
\textbf{Author/Year} & \textbf{Model Formulation} & \textbf{Dataset (System, $\sqrt{s}$)} & \textbf{$p_T$ Range [GeV/c]} & \textbf{Typical $\chi^2$/ndf} & \textbf{Key Physical Finding} \\ \hline
Schnedermann (1993) & BGBW (Hydrodynamic) & SPS (Pb-Pb, 200 AGeV) & $0 - 2.5$ & $\sim 1.0 - 2.0$ & Introduced collective radial flow profile. \\ \hline
Cleymans (2012) & Thermodynamic Tsallis & ALICE ($pp$, 0.9, 2.76, 7 TeV) & $0.5 - 10$ & $\sim 0.5 - 1.5$ & Derived consistent Tsallis thermodynamics. \\ \hline
Sahoo (2015) & Tsallis-Pareto & CMS ($pp$, 0.9, 2.76, 7 TeV) & $0.1 - 200$ & $< 1.5$ & Described full spectrum up to extreme $p_T$. \\ \hline
Tang (2009) & Tsallis Blast-Wave (TBW) & RHIC (Au-Au, 200 GeV) & $0 - 6$ & $\sim 1.0$ & Embedded non-extensivity in flow field. \\ \hline
Biro (2020) & Tsallis-Thermometer & ALICE ($pp$, 7 TeV) & $0 - 20$ & $\sim 1.2$ & Connected $q$ to temperature fluctuations. \\ \hline
Khuntia (2019) & BGBW + Tsallis & ALICE ($pp$, 7 TeV) & $0 - 15$ & $\sim 1.0 - 2.5$ & Mass ordering in $T_{\text{kin}}$, flat $\langle \beta \rangle$ vs mult. \\ \hline
Rath (2020) & BGBW + Tsallis & ALICE ($pp$, 5.02, 13 TeV) & $0 - 20$ & $\sim 0.8 - 2.0$ & Energy independence of radial flow. \\ \hline
Grigoryan (2017) & Modified Tsallis & ATLAS ($pp$, 0.9, 7 TeV) & $0.5 - 50$ & $\sim 1.5$ & Scaling behavior of Tsallis parameters. \\ \hline
Biro (2025) & 2-Comp Soft/Hard & ALICE ($pp$, 2.76 - 13 TeV) & $0 - 20$ & $< 1.0$ & Hard residual invariant with multiplicity. \\ \hline
Wong (2015) & Tsallis with Hard Scatt. & CMS/ALICE ($pp$, 7 TeV) & $1 - 100$ & $\sim 1.0$ & Linking $q$ parameter to QCD calculus. \\ \hline
\end{tabular}
}
\end{table}

\section{AI Methods in Scientific Validation}
The volume and complexity of scientific literature have catalyzed the adoption of Artificial Intelligence (AI) and Large Language Models (LLMs) in the research process. A critical development is Retrieval-Augmented Generation (RAG) \citep{Lewis2020}, which grounds LLM outputs in verified external corpora, mitigating hallucinatory claims. In physics, RAG systems are increasingly utilized to navigate the arXiv preprint repository, extracting specific experimental parameters or theoretical assumptions to contextualize new results.

Beyond literature retrieval, LLMs are being actively investigated for automated peer review and theorem proving. \citet{Wang2023} demonstrated the utility of LLMs in identifying methodological flaws in scientific manuscripts, while \citet{Duffield2024} explored multi-agent frameworks (such as the ReAct paradigm) for autonomous theorem verification in formal mathematics. In phenomenological physics, bridging the gap between symbolic derivation (e.g., via SymPy) and numerical execution (e.g., MINUIT) using AI orchestration represents a frontier application. This thesis leverages these advancements, employing an MCP (Model Context Protocol) toolchain to grant AI agents direct access to symbolic math evaluators and numerical optimizers.

\section{Reproducibility and Model Selection Criteria}
The reproducibility crisis emphasizes the need for objective mathematical criteria when comparing competing phenomenological models, particularly nested or multi-component formulations. Minimizing $\chi^2$/ndf is necessary but insufficient, as adding parameters inevitably improves the fit residual. 

To penalize over-parameterization, Information Criteria are employed. The Akaike Information Criterion (AIC) and Bayesian Information Criterion (BIC) are defined as:
\begin{align}
\text{AIC} &= 2k - 2\ln(\hat{L}) \\
\text{BIC} &= k\ln(n) - 2\ln(\hat{L})
\end{align}
where $k$ is the number of free parameters, $\hat{L}$ is the maximized value of the likelihood function, and $n$ is the number of data points. For models fitted via least squares with normally distributed errors, $-2\ln(\hat{L})$ can be approximated by $\chi^2$. Furthermore, the F-test provides a frequentist method to determine if the reduction in $\chi^2$ achieved by adding a new parameter component is statistically significant relative to the loss of degrees of freedom. Incorporating these criteria into an automated validation pipeline is essential for distinguishing physical discovery from mathematical overfitting.
